\section{Introduction}
\label{sec:intro}

PolicyEngine Macro is a suite of open-source macroeconomic models built around a single reform vocabulary: a tax--benefit reform is a flat dictionary of PolicyEngine parameter changes, and every member of the suite consumes that same object through its own declared contract. The suite's other members are anchored to published figures --- an emulator of the Office for Budget Responsibility's (OBR) macroeconometric model, a Bank-of-England-style structural VAR, and the PolicyEngine static microsimulation itself. This paper documents the member that answers a different kind of question: the \emph{psl-og} model, a dynamic overlapping-generations (OLG) general-equilibrium model built on the Policy Simulation Library's OG-Core framework \citep{DeBackerEvans2024}, wired into the suite today through its United Kingdom country calibration OG-UK \citep{OGUKdocs}, with the United States calibration OG-USA \citep{DeBackerEvansOGUSA} in scope.

The model is a large-scale descendant of \citet{AuerbachKotlikoff1987}. It tracks eighty annual age cohorts --- households aged 20 to 99, entering economic life at 20 and facing age-specific mortality --- each populated by several lifetime-ability types. Every cohort chooses labour supply and saving over its remaining lifetime, trading off consumption against an elliptical disutility of work and a warm-glow bequest motive. Perfectly competitive firms rent the resulting capital and labour; the government taxes, transfers, spends, and issues debt, with a fiscal rule that stabilises the long-run debt-to-GDP ratio. Two solution objects are on offer: the \emph{stationary steady state} (``where does the economy settle under this policy?'') and the \emph{transition path} (a roughly 60-year, year-by-year solve of how it gets there).

Within the suite, the OLG member's role is long-run structural analysis. The OBR emulator and the SVAR speak to demand-side dynamics over a forecast window; the static microsimulation speaks to first-round revenue and distribution. The OLG model is the only member in which a tax reform changes prices --- wages, interest rates --- through households' labour-supply and saving responses, so it is the member that can address questions like the long-run output cost of a higher basic rate of income tax, the crowding-out effects of debt, or the incidence of corporation tax across generations. Its distinguishing input channel is the microdata bridge of \citet{DeBackerEvansPhillips2019}: reforms are not stylised wedges but actual statute, run through the PolicyEngine UK microsimulation model \citep{PolicyEngineUK} over enhanced Family Resources Survey (FRS) microdata and compressed into smooth parametric effective- and marginal-tax-rate functions that the optimising households face.

Two properties of this member must be stated at the outset, because they shape everything downstream. First, it is \emph{local-only}: a single steady-state solve --- microdata calibration plus the general-equilibrium fixed point --- has been measured at over seventeen minutes on a laptop, and a full transition path takes hours, so the model is deliberately excluded from the hosted PolicyEngine Macro service and runs only on the user's own machine through the \texttt{pe-macro} command-line interface. Second, unlike the suite's OBR and SVAR members, it is a \emph{calibrated structural model with no published ground truth to replicate}. Its outputs are structural counterfactuals --- internally consistent answers conditional on the model's assumptions --- not validated forecasts. This paper is accordingly written as documentation with explicit caveats (Section~\ref{sec:limitations}), not as a validation exercise.

The remainder of the paper proceeds as follows. Section~\ref{sec:literature} places the model in the OLG literature and cites the framework's own documentation. Section~\ref{sec:model} sets out the model's equations: the household problem, the firm problem, the government budget and fiscal rules, market clearing, and the steady-state and transition-path solution algorithms. Section~\ref{sec:taxfunc} details the tax-function estimation pipeline from PolicyEngine microdata --- the PolicyEngine-specific contribution. Section~\ref{sec:implementation} describes the software architecture and the \texttt{policyengine-macro} integration. Section~\ref{sec:calibration} documents the UK calibration and compares selected targets with official statistics. Section~\ref{sec:limitations} states the model's limitations honestly, and Section~\ref{sec:conclusion} concludes. Three appendices supply a parameter glossary, the OG-Core API surface the integration touches, and an exact pinned-version reproduction recipe. The figures throughout are generated from the calibration's cheap-to-load inputs --- the UN demographic series, the OG-Core ability matrix, the fitted utility and tax functional forms --- and require no equilibrium solve to reproduce (\texttt{figures/make\_figures.py}).
